\documentclass[10pt,a4paper]{article}
\usepackage[T1]{fontenc}
\usepackage[margin=21mm]{geometry}
\usepackage{lmodern,amsmath,amssymb,amsthm,mathtools,microtype}
\usepackage[hidelinks]{hyperref}
\newtheorem{theorem}{Theorem}
\title{The unit in a relative Erd\H{o}s--Straus boundary packet}
\author{The Clankers}
\date{16 September 2026 --- research preprint}
\begin{document}\maketitle
\begin{abstract}
An available positive divisor packet may reduce to a subgroup norm times a
non-monomial unit. Retaining this unit gives an integer lower bound and an
original-word inverse. An explicit all-precision family has one middle cofactor,
requiring $2^{k-5}+1$ prime occurrences in either role, while the preceding
canonical one-/two-prime block tests fail. Three members have deterministic
primality certificates. No universal ES or historical-priority claim is made.
\end{abstract}

\paragraph{Original arithmetic.}
Fix $k\ge3$, $u=2^{2k-5}$, $p\equiv1\pmod8$, $p>2u$, and $N=p+4u$.
The original middle states at this seed correspond to
\[
 Q\mid N,\quad Q\equiv-1\pmod{2^k},\quad R=N/Q,\quad a=(p+R)/4.
 \tag{1}
\]
Indeed $u\mid a^2\iff2^{k-2}\mid a$, and (1) gives that congruence.
Both factors are $3\pmod4$ and at least three; $N<3p$ makes them less than $p$.
Also $4(a+u)=R+N$ gives $R\mid a+u$. Conversely those conditions imply (1).
With $g=\gcd(a,u)$, valuations give
\[
 (h,r,s)=(g^2/u,u/g,a/g),\quad a=hrs,\quad u=hr^2,\quad\gcd(r,s)=1.
\]
All factors are units modulo $R$, so $\lambda=(r+s)/R$ is integral and
\[
 (x,y,z)=(a,phs\lambda,phr\lambda),\quad
 \frac4p=\frac1x+\frac1y+\frac1z,\quad
 u=\frac{pa^2}{Ry-pa}.                                    \tag{2}
\]
The reciprocal identity follows by multiplying by $phrs\lambda$.
These are inherited Type II coordinates \cite{ET}. Structural results allow
integer $p$; primality is a separately certified input in examples.

\begin{theorem}[Unit-corrected positive lift]\label{lift}
Let $o$ be a power of two, $1<D<o$ a divisor, and
$P\in\mathbb Z_{\ge0}[C_o]$ an actual labelled packet count. In
$A_o=\mathbb F_2[X]/(X^o-1)=\mathbb F_2[T]/T^o$, $T=X+1$, suppose
\[
 \overline P=T^\nu U(T),\quad U(0)=1,\quad o-D\le\nu<o.
\]
Put $q_T=T^{\nu-o+D}U\bmod T^D$, $q(X)=q_T(X+1)\in A_D$, and let
$\widehat q$ be its $0/1$ coefficient lift. There is a unique
$E\in\mathbb Z_{\ge0}[C_o]$ with
\[
 \boxed{P=\mathcal N_D\widehat q+2E,\qquad
 \mathcal N_D=\sum_{j=0}^{o/D-1}X^{jD}=T^{o-D}.}             \tag{3}
\]
For every nonnegative original complementary packet $G$,
\[
 \boxed{[X^t]PG\ge[X^{t\bmod D}](\widehat q\,\pi_DG),}    \tag{4}
\]
where $\pi_D$ sums whole integer coefficients over the congruence classes.
If $\nu=o-D$, the module map $h\mapsto\mathcal N_Dqh$ identifies $A_D$ with
$\ker T^D$. Its inverse first reads constant coefficients on $D$-cosets,
then multiplies by $q^{-1}$. Ordinary quotient summation is not the inverse.
\end{theorem}
\begin{proof}
Multiplication by the norm repeats each coefficient on a coset. Thus the
parity of $P$ is $q$ repeated. Subtracting that $0/1$ vector gives the unique
nonnegative even remainder (3). Convolution gives the stronger exact identity
$[X^t]PG=[X^{t\bmod D}]\widehat q\pi_DG+2[X^t]EG$.
The $T$-coordinate map $h\mapsto T^{o-D}h$ has image $\ker T^D$ and is injective
on $A_D$. When $q_T(0)=1$, its inverse is
$\sum_{j=0}^{D-1}(1+q_T)^j$. This proves all assertions.
\end{proof}
For $z=\nu-o+D>0$, the same map has kernel $T^{D-z}A_D$ and image $T^\nu A_o$;
these exceptions are not inverted. The boundary source is SplitZero BR1 with
prequotient $\mathbb F_2[[T]]$, $v=T^D$, $f=T^{o-D}U$ \cite{SZ}.
Both multipliers are injective before quotienting, and $U$ is a power-series
unit. The original supported zero $e=0_A^\bullet$ is not absence $\tau$:
even positive coefficients and their original preimages remain in (3).

\newpage
\paragraph{Constructive input and original-word inverse.}
In the cyclic chart $H_b=\langle b\rangle$, $b=3$ or $-3$, retain a packet
$W=A\prod_iG_i^{t_i}$ on disjoint original prime blocks, with affine exponents
$v_q(W)=a_q+\sigma_qt_i$ and $0\le t_i\le E_i$. Endpoint inequalities put
every exponent in its actual prime-factor inventory. Require $A\notin H_b$,
$G_i\in H_b$; retain the formal role $p^{-\varepsilon}$, $\varepsilon=0,1$.
This role is not a prime factor. Let the cyclic step logs be $a_i$, of order
$o=2^{k-2}$.

Choose actual subcapacities $c_i\le E_i$, put
$\nu=\sum_i c_i2^{v_2(a_i)}<o$, and use
\[
 P(X)=\prod_i\sum_{t_i\preceq c_i}X^{a_it_i},\qquad
 U=\prod_i\left(\frac{1+(1+T)^{a_i}}{T^{2^{v_2(a_i)}}}\right)^{c_i}. \tag{5}
\]
Here $t\preceq c$ is binary-submask containment. Frobenius proves
$\overline P=T^\nu U$ without replacing repeated occurrences by coloured
copies. Untouched coordinates can supply $G$; a used coordinate is not also
an independent remaining factor.

A positive lower coefficient in (4) provides an actual $G$-word and an odd
fine coefficient of $P$. Factor (5) into its binary factors and compute suffix
parities. At an odd target, exactly one of the two next suffix coefficients
is odd; following it terminates at exponent zero. The returned $t_i$ are
actual submasks. Affine inversion reads
$(v_q(W)-a_q)/\sigma_q$ at a nonzero step and checks all other coordinates.
The final equation $Wp^{-\varepsilon}=-1\pmod{2^k}$ returns $Q=W$ at role zero
or $Q=N/W$ at role one, then (1)--(2). Fixed-role words are distinct; forgetting
the role identifies at most the two complementary words. Thus (4) bounds role
words, or half that many original states rounded upwards when both roles occur.

\paragraph{A nontrivial integral boundary as well.}
Let $D$ be a power of two and $h<D$ odd. For $q=X^r(1+\cdots+X^{h-1})$, write
$a=h^{-1}\bmod D$ in $[1,D)$ and $b=(ha-1)/D$. Direct geometric multiplication gives
\[
 q^{-1}=X^{-r}\left(\sum_{j=0}^{a-1}X^{hj}-\frac bh\sum_{j=0}^{D-1}X^j\right)
 \quad\text{in }\mathbb Q[C_D].                            \tag{6}
\]
Therefore its integer circulant has image exactly the vectors whose coefficient
sum is divisible by $h$: necessity is its column sum; sufficiency follows from
(6). Its cokernel is $\mathbb Z/h$ and determinant has absolute value $h$.
Its cyclic coefficient Gram is $(o/D)C_q^{\mathsf T}C_q$, not the identity.
In particular $q=X^3+\cdots+X^7\in A_8$ is invertible modulo two but has integer
index five. The inverse has signs and denominators; it is not used as a
positivity-preserving inverse. The original-word quotient instead weights an occupied
residue by $1/P_r$; a formal periodic amplitude can lift only if it vanishes
where $P_r=0$. These source constraints are retained in the full proof.

\begin{theorem}[A unique-state family with unbounded required length]
Let $k\ge6$, $d=2^{k-4}$, and choose distinct primes $q,r,s\ne3$ with
\[
 q\equiv3^{3d-1},\qquad r\equiv-3^{5d/2},\qquad
 s\equiv-3^{d/2+2}\pmod{2^k}.
\]
Set $N=3^{d-1}qrs$, $p=N-4u$, requiring $p>2u$. Then $p\equiv1\pmod{24}$,
and the complete fixed-seed branch has exactly
\[
 \boxed{Q=3^{d/2-1}qs,
 \qquad R=3^{d/2}r,
 \qquad \Omega(Q)=\Omega(R)=2^{k-5}+1.}                     \tag{7}
\]
Theorem~\ref{lift} forces that state using a proper boundary and a non-monomial
unit. The preceding canonical disjoint singleton/two-prime full and relative
packet tests fail throughout this family.
\end{theorem}
\begin{proof}
In $H_3$, only $r,s$ are outside, each once. Inside words have logs in
$[0,d-1]\cup[3d-1,4d-2]$. The $r$ target $3d/2$ misses both intervals;
the $s$ target $7d/2-2$ has the unique expression $(3d-1)+(d/2-1)$.
This proves (7). The phase is $3d$, so the role has log $d$. The actual
constant-$s$ packet satisfies the \emph{integer} identity
\[
 (1+\cdots+X^{d-1})(1+X^{3d-1})(1+X^d)
 =\mathcal N_{2d}(X^{d-1}+\cdots+X^{2d-1})+2(1+\cdots+X^{d-2}). \tag{8}
\]
The target reduces to $3d/2-2$, an occupied coefficient of the displayed unit.
This proves the new certificate, with no target-availability assumption.
\end{proof}

\newpage
\paragraph{Exact scope of the separation.}
For the preceding test class, let $D_0$ be its common dyadic lower stride.
A step log $a$ of capacity $E$ has high stride
$\delta=D_0/\gcd(a,D_0)$ and high weighted capacity at most
$\lfloor E/\delta\rfloor\operatorname{lowbit}(a\delta/D_0\bmod(o/D_0))$.
The complete upper bounds for all its disjoint singleton/same-coset pair
patterns, including the role, are
\[
\begin{array}{c|cccc}
D_0&1&2&4\le D_0\le d&2d\\\hline
\text{maximum upper bound}&2d+2&d+1&2d/D_0&0.
\end{array}                                                 \tag{9}
\]
Each is below the required $4d/D_0-1$. To verify the list for all $d$, in the
$3$-chart the two inside singletons have weights/capacities $(1,d-1),(1,1)$.
Their pair has weight $d$ or two and capacity one. The outside pair has weight
two when active; outside singletons have zero capacity. The role has weight
$d$ and capacity one. At stride at least four the small capacity-one steps
vanish; only the role, the weight-$d$ pair, or at most $d/D_0-1$ singleton
units remain. This gives (9), with strides one and two calculated directly.

In the $-3$-chart all four primes are outside. Only the $3$ singleton can have
positive capacity, namely $d/2-1$ at weight two. The $(3,q)$ pair has weight
$d$ or two; each cross pair has weight one; the $(r,s)$ pair has weight two.
Every pair has capacity at most one. Disjointness and the same role again give
(9). Lower digits only reduce these bounds. This proves failure before the old
quotient target is considered.

A different map does survive: the \emph{three-prime shear}
\[
(v_3,v_r,v_s)=(a+2t,t,1-t),\quad0\le a\le d-3,\quad t=0,1
\]
has half-turn step $3^2r/s=3^{2d}\pmod{2^k}$ and inverse
$t=v_r$, $a=v_3-2v_r$. With $q$ retained and $a=d/2-1$, its $t=0$ endpoint is
(7). It is not a disjoint pair block. No failure of all higher-arity maps is claimed.

Dirichlet's theorem supplies infinitely many auxiliary choices. The derived
$p$ are not asserted prime. They can all be $1\pmod{840}$ by additionally
requiring $q\equiv(1+4u)(3^{d-1}rs)^{-1}\pmod{35}$, with $r,s$ avoiding $5,7$.
The residue is a unit since $1+2^{2k-3}$ is nonzero modulo $5,7$ \cite{Dirichlet}.

\paragraph{Certified prime members.}
The auxiliary choices $(6,59,23,47)$, $(7,17099,239,167)$ and
$(8,1259,223,599)$ for $(k,q,r,s)$ give the primes
\[
1721521,\quad1492567108321,\quad2413105093468609.
\]
Their unique cofactor/residual pairs are
$(8319,207)$, $(77099391,19359)$ and $(1649306367,1463103)$, with lengths
three, five and nine in \emph{both} roles. Complete-order primality certificates
and all divisor enumerations accompany the source. For the first prime,
\[
\frac4{1721521}=\frac1{430432}+\frac1{12041213064920}+\frac1{3580763680}.
\]
For the largest prime the state is
\[
(u,a,h,r,s,\lambda)=(2048,603276273732928,2,32,9426191777077,6442603).
\]
The finite comparison through two million adds only one prime to the prior
union of relative, two-prime and $3,11$ tests. A three-occurrence search with
prior tests already covers every occupied branch there. The unbounded family,
not an inflated small-range comparison, establishes the structural separation.
No universal ES theorem, new global verification range, Lean build, independent
mathematical review, or historical-priority conclusion is asserted.

\begingroup\footnotesize
\begin{thebibliography}{4}
\bibitem{ET} Elsholtz, C., \& Tao, T. (2013). Counting the number of solutions
to the Erd\H{o}s--Straus equation on unit fractions.
\emph{J. Aust. Math. Soc., 94}(1), 50--105. arXiv:1107.1010, \S2, (2.13)--(2.22).
\bibitem{SZ} The Clankers (2026). \emph{Mixed-boundary recovery on the original
SplitZero quotient}, BR1--BR2. Zeta workbench commit
\texttt{1efd53337561d8f67cf0ac1d119acdaa222644c3}.
\bibitem{Dirichlet} Sutherland, A. V. (2017). \emph{Dirichlet L-functions,
primes in arithmetic progressions}. MIT 18.785 Lecture 18, Theorem 18.1, p.~1.
\bibitem{MCJ} Massey, J. L., Costello, D. J., Jr., \& Justesen, J. (1973).
Polynomial weights and code constructions. \emph{IEEE Trans. Inf. Theory, 19}(1),
101--110. Classical weight-retention antecedent; the full note gives a direct
binary proof and the exact scope of the source reading.
\end{thebibliography}\endgroup
\end{document}
